        Das Programm, um 5 Zufallszahlen auszugeben, sieht zum Beispiel folgendermaßen aus:
\begin{code}
#include <stdio.h>
#include <stdlib.h>

int main() {
    for(int i = 1; i <= 5; i++) {
        int zahl = rand();
        printf("Zufallszahl %d: %d\n", i, zahl);
    }

    return 0;
}
\end{code}
dessen Ausgabe folgende ist:
\begin{verbatim}
Zufallszahl 1: 1804289383
Zufallszahl 2: 846930886
Zufallszahl 3: 1681692777
Zufallszahl 4: 1714636915
Zufallszahl 5: 1957747793
\end{verbatim}
    Sind diese Zahlen die Gleichen, die auch Ihr seht?
    (Ihr könnt nun im Text nach der Aufgabe im Kapitel~\ref{sec:c:wuerfel} weiter machen, dort findet Ihr die Erklärung für dieses Verhalten.)
